% Prospective assistance/retention protocol; no biological experiment reported.
\section{From a guided action to a retained skill}

Catch a ball with outside guidance, then try again after the guidance ends. The
second catch asks what changed in the trainee. The comparison is useful only if
the task and access to movement remain comparable; removing a necessary movement
bridge along with the trainer would confound the test. This distinction carries
the proposed sequence from touch and object handling to NeuroBoxer and NeuroArcher.

The near-term apparatus would record scalp EEG and surface EMG from sleeves on
both arms, alongside synchronized, calibrated left and right tripod views.
OpenBCI or an equivalent configurable system remains a candidate requiring
channel, site, reference, bandwidth, unit, delay and clock qualification. Each
event needs its independently observed movement and, where claimed, an
appropriate contact or force reference. Surface EMG supplies neither direct
force nor afferent nerve activity. A virtual body can predict a response, but
its predicted contact or sensory state cannot stand in for those observations.

This left/right hand-task apparatus and the earlier opposing top/bottom dot-map
apparatus are different proposed configurations of the motion-capture
relay-composite. Their calibration, frames, visibility, occlusion and timing
must be qualified separately. Preserve the source of each measured point and
the uncertainty of the joined map. The musculoskeletal-simulation physics-engine
also needs task- and organism-specific anatomy and mechanics; lifelike rendering
cannot establish its function.

\subsection{Build coordination through four declared tasks}

\paragraph{Touch and movement.}
Rapid hand rubbing brings contact and movement together. Record the requested
knuckle movement and interlaced-hand stretch, with palms outward and arms
extended, as separate conditions. Matched observations must distinguish
contact-related effects from movement-related effects; the cold-hands metaphor
does not establish a cold-exposure result or identify a unique sensation.

\paragraph{Ball transfer and juggling.}
The proposed hot-potato task passes one juggling ball between the hands, followed
by the preferred two-ball juggling task. One limb's action changes the next
limb's problem. Independently observed release and catch events, trajectories
and timing would test that coordination. Preserve uncertainty when an event or
body part is not observed, and compare with cue- and history-based controls.

\paragraph{Punch discrimination and NeuroBoxer.}
Independently labelled punch classes would test discrimination and movement
prediction before becoming the NeuroBoxer virtual game. Boxing includes hand
and wrist stabilization as well as arm motion. A prime-era Mike Tyson opponent
is a proposed virtual game character. A virtual opponent supplies no physical
opponent or evidence of safety. Rights-qualified expert footage can inform
movement and tactics, but supplies no expert neural, EMG or force recordings.

\paragraph{Toy archery and NeuroArcher.}
Toy archery would add explicit grasp, asymmetric two-hand roles, draw under
opposing loads, release and return. Its proposed arrows have small, cup-shaped
suction tips and are carried in an over-arm pouch with multiple arrows. NeuroArcher would use that coordination to act
on a virtual target. Success still requires a declared outcome and an
independent comparison, rather than agreement between simulated components.

The whole snap is a further coordinated showcase, from arm-out/palm-down rest
through elbow and forearm swing and turn, wrist and hand upward during the snap,
then downward return. The cup reach remains a simpler precision example. The
motion, contact and anatomy needed for either task are specific; a human-looking
animation cannot demonstrate a monkey's ability to perform the snap.

\subsection{Separate motion control, neural contribution and physical writing}

The first game milestone uses body motion to control a virtual body. A later
EEG-only milestone would ask whether permitted brain observations contribute to
the specified action. Withheld EMG and measured movement cannot become predictor
inputs through alignment, filtering or cue leakage. Eye activity, covert muscle
activity, residual movement and shared task information need their own controls.
Keep the original brain-only, body-only and paired comparisons, target-persistence
baseline, partition rules and success criteria.

Aim to make the declared intent-to-action mapping as close to one-to-one as
possible by measuring its intermediary. State the participant, session, condition,
allowed context, tolerated error and independent reference. Reduced ambiguity
would be a task-specific result; it would not identify a unique hidden intention
or establish the best achievable mapping. Reverse reconstruction of recorded
EEG features remains conditional and cannot deliver sensation.

The proposed later physical writer is a separate experiment. EEG and surface
EMG record, whereas neuromuscular electrical stimulation (NMES) and functional electrical stimulation (FES) denote stimulation approaches that can produce
contractions or functional movement. Sensation delivery needs its own writer or
sensory transducer. Device, clinical, consent and physical admission remain
separate, and no programme spinal intervention or qualified spinal-writing interface is reported. These writers and
actuators are unqualified and outside the existing acquisition/virtual-work
budget scope; no recording or intervention is reported here.

\subsection{Ask what remains after the help ends}

A proposed blindfolded condition changes visual feedback. It cannot independently
establish neural contribution, causal control or retained skill. Those claims
need matched comparisons, including a trainer-off test that preserves any movement
bridge needed to perform the task. Measure immediate performance, delayed
retention and transfer to specified unfamiliar conditions separately. Report
remaining support, uncertainty and alternative explanations.

The model's training and the trainee's learning also need separate evidence.
Qualified paired data could fit predictions; optional control-policy reinforcement
learning needs states, actions, transitions and rewards. A simulated reward is
not evidence of biological teaching. Final held-out outcomes remain outside
fitting, fine-tuning, rewards and model selection; frozen inference is distinct
from a separately reviewed adaptive update. The retained LINK comparison failed
both model-admission gates and all five benefit claims. These prospective tasks
do not revise that result or establish retained learning, tissue repair or a cure.
